% =====================================================================
% QR-CODE: Excel-Spielwiese
% Informatik Klasse 7 | Tabellenkalkulation Stunde 2
% =====================================================================

\documentclass[a4paper,11pt]{article}

\usepackage[utf8]{inputenc}
\usepackage[T1]{fontenc}
\usepackage{helvet}
\renewcommand{\familydefault}{\sfdefault}

\usepackage[margin=20mm]{geometry}
\usepackage{tikz}
\usepackage{xcolor}
\usepackage{qrcode}
\usepackage{hyperref}
\hypersetup{colorlinks=false, pdfborder={0 0 0}}

% =====================================================================
% FARBEN
% =====================================================================

\definecolor{informatik}{HTML}{b35c00}
\colorlet{fachfarbe}{informatik}

% =====================================================================
% KONFIGURATION
% =====================================================================

\newcommand{\simtitel}{Excel-Spielwiese}
\newcommand{\simurl}{https://devbox42.github.io/sciencesim/informatik/kl07/02-tabellenkalkulation/SIM-02-excel-spielwiese.html}
\newcommand{\simurlkurz}{devbox42.github.io/.../SIM-02-spielwiese.html}

% =====================================================================
% DOKUMENT
% =====================================================================

\pagestyle{empty}

\begin{document}

\begin{center}

% Titel
\vspace*{10mm}
{\fontsize{28}{34}\selectfont\bfseries\color{fachfarbe}\simtitel}

\vspace{15mm}

% QR-Code mit Rahmen (anklickbar)
\href{\simurl}{\fcolorbox{fachfarbe}{white}{\qrcode[height=100mm]{\simurl}}}

\vspace{12mm}

% URL
{\large\color{gray}\simurlkurz}

\vspace{8mm}

% Hinweis
{\Large\color{gray}Scanne den QR-Code mit deinem Gerat}

\vspace{10mm}

% Info-Box
\fcolorbox{fachfarbe}{fachfarbe!10}{%
    \parbox{0.75\textwidth}{%
        \centering
        \vspace{3mm}
        Experimentiere frei mit Formeln und Funktionen!\\[2mm]
        Lade verschiedene Beispiel-Daten und probiere alle Funktionen aus.
        \vspace{3mm}
    }%
}

\end{center}

\end{document}
